\documentclass[10pt,a4paper]{amsart}
\usepackage{amsmath,amssymb,amsthm,graphicx}
\usepackage{verbatim}
\newenvironment{reference}{\comment}{\endcomment}
%\newenvironment{reference}{}{}
\newenvironment{slogan}{\comment}{\endcomment}
\newenvironment{history}{\comment}{\endcomment}

% For commutative diagrams we use Xy-pic
\usepackage[all]{xy}

% We use 2cell for 2-commutative diagrams.
\xyoption{2cell}
% Theorem environments.
%
\theoremstyle{plain}
\newtheorem{theorem}[subsection]{Theorem}
\newtheorem{proposition}[subsection]{Proposition}
\newtheorem{lemma}[subsection]{Lemma}

\theoremstyle{definition}
\newtheorem{definition}[subsection]{Definition}
\newtheorem{example}[subsection]{Example}
\newtheorem{exercise}[subsection]{Exercise}
\newtheorem{situation}[subsection]{Situation}

\theoremstyle{remark}
\newtheorem{remark}[subsection]{Remark}
\newtheorem{remarks}[subsection]{Remarks}

\numberwithin{equation}{subsection}

% Macros
%
\def\lim{\mathop{\mathrm{lim}}\nolimits}
\def\colim{\mathop{\mathrm{colim}}\nolimits}
\def\Spec{\mathop{\mathrm{Spec}}}
\def\Hom{\mathop{\mathrm{Hom}}\nolimits}
\def\Ext{\mathop{\mathrm{Ext}}\nolimits}
\def\SheafHom{\mathop{\mathcal{H}\!\mathit{om}}\nolimits}
\def\SheafExt{\mathop{\mathcal{E}\!\mathit{xt}}\nolimits}
\def\Sch{\mathit{Sch}}
\def\Mor{\mathop{\mathrm{Mor}}\nolimits}
\def\Ob{\mathop{\mathrm{Ob}}\nolimits}
\def\Sh{\mathop{\mathit{Sh}}\nolimits}
\def\NL{\mathop{N\!L}\nolimits}
\def\CH{\mathop{\mathrm{CH}}\nolimits}
\def\proetale{{pro\text{-}\acute{e}tale}}
\def\etale{{\acute{e}tale}}
\def\QCoh{\mathit{QCoh}}
\def\Ker{\mathop{\mathrm{Ker}}}
\def\Im{\mathop{\mathrm{Im}}}
\def\Coker{\mathop{\mathrm{Coker}}}
\def\Coim{\mathop{\mathrm{Coim}}}

% Boxtimes
%
\DeclareMathSymbol{\boxtimes}{\mathbin}{AMSa}{"02}

%
% Macros for moduli stacks/spaces
%
\def\QCohstack{\mathcal{QC}\!\mathit{oh}}
\def\Cohstack{\mathcal{C}\!\mathit{oh}}
\def\Spacesstack{\mathcal{S}\!\mathit{paces}}
\def\Quotfunctor{\mathrm{Quot}}
\def\Hilbfunctor{\mathrm{Hilb}}
\def\Curvesstack{\mathcal{C}\!\mathit{urves}}
\def\Polarizedstack{\mathcal{P}\!\mathit{olarized}}
\def\Complexesstack{\mathcal{C}\!\mathit{omplexes}}
% \Pic is the operator that assigns to X its picard group, usage \Pic(X)
% \Picardstack_{X/B} denotes the Picard stack of X over B
% \Picardfunctor_{X/B} denotes the Picard functor of X over B
\def\Pic{\mathop{\mathrm{Pic}}\nolimits}
\def\Picardstack{\mathcal{P}\!\mathit{ic}}
\def\Picardfunctor{\mathrm{Pic}}
\def\Deformationcategory{\mathcal{D}\!\mathit{ef}}
\usepackage{cleveref}
\crefname{originallemma}{original lemma}{original lemmas}
\Crefname{originallemma}{Original lemma}{Original lemmas}
\makeatletter
\newlabel{lemma-quotient-stack-arrows}{{044R}{source}{}{sourcepointer-1}{}}
\newlabel{lemma-quotient-stack-arrows@cref}{{[originallemma][0][]044R}{[1][0][]source}{}{}{}}
\makeatother
\crefname{originallemma}{original lemma}{original lemmas}
\Crefname{originallemma}{Original lemma}{Original lemmas}
\makeatletter
\newlabel{lemma-quotient-stack-2-arrow}{{044S}{source}{}{sourcepointer-2}{}}
\newlabel{lemma-quotient-stack-2-arrow@cref}{{[originallemma][0][]044S}{[1][0][]source}{}{}{}}
\makeatother
\crefname{originalreference}{original reference}{original references}
\Crefname{originalreference}{Original reference}{Original references}
\makeatletter
\newlabel{stacks-lemma-2-product-stacks-in-groupoids}{{02ZL}{source}{}{sourcepointer-3}{}}
\newlabel{stacks-lemma-2-product-stacks-in-groupoids@cref}{{[originalreference][0][]02ZL}{[1][0][]source}{}{}{}}
\makeatother
\crefname{originalreference}{original reference}{original references}
\Crefname{originalreference}{Original reference}{Original references}
\makeatletter
\newlabel{categories-lemma-2-product-categories-over-C}{{0040}{source}{}{sourcepointer-4}{}}
\newlabel{categories-lemma-2-product-categories-over-C@cref}{{[originalreference][0][]0040}{[1][0][]source}{}{}{}}
\makeatother
\crefname{originallemma}{original lemma}{original lemmas}
\Crefname{originallemma}{Original lemma}{Original lemmas}
\makeatletter
\newlabel{lemma-quotient-stack-morphisms}{{044V}{source}{}{sourcepointer-5}{}}
\newlabel{lemma-quotient-stack-morphisms@cref}{{[originallemma][0][]044V}{[1][0][]source}{}{}{}}
\makeatother
\crefname{originallemma}{original lemma}{original lemmas}
\Crefname{originallemma}{Original lemma}{Original lemmas}
\makeatletter
\newlabel{lemma-cartesian-square-of-morphism}{{04Y4}{source}{}{sourcepointer-6}{}}
\newlabel{lemma-cartesian-square-of-morphism@cref}{{[originallemma][0][]04Y4}{[1][0][]source}{}{}{}}
\makeatother
\title{The quotient-stack square is a 2-fibre product of stacks in groupoids}
\author{The Stacks Project Authors}
\date{stacks-project:spaces-groupoids.tex:local-snapshot-20261012; GNU FDL 1.2 or later}
\begin{document}
\maketitle
\noindent\textit{Editorial note.} All delivered mathematical text is the author's own, copied verbatim from the pinned Stacks Project chapter file spaces-groupoids.tex; no statement or proof step was rewritten, reordered or removed.\par\medskip
\noindent\textit{Editorial note.} Span recomputation by label pairing (LANE\_BRIEF 10.2): in this chapter the author's proof environment follows the statement environment as a standalone block; the spec covers the complete environment pair (statement span, proof span) as recorded in the pool candidate.\par\medskip
\ifdefined\hypertarget\hypertarget{sourcepointer-1}{}\fi
\noindent\textit{Source reference.} Original lemma 044R: sources/stacks-chow/tags.txt, source lines 5410--5410. Cross-chapter standard prerequisite cited by the author inside the counted span; it is used as a pointer and its own proof is not part of this item.\par\medskip
\ifdefined\hypertarget\hypertarget{sourcepointer-2}{}\fi
\noindent\textit{Source reference.} Original lemma 044S: sources/stacks-chow/tags.txt, source lines 5411--5411. Cross-chapter standard prerequisite cited by the author inside the counted span; it is used as a pointer and its own proof is not part of this item.\par\medskip
\ifdefined\hypertarget\hypertarget{sourcepointer-3}{}\fi
\noindent\textit{Source reference.} Original reference 02ZL: sources/stacks-chow/tags.txt, source lines 3916--3916. Cross-chapter standard prerequisite cited by the author inside the counted span; it is used as a pointer and its own proof is not part of this item.\par\medskip
\ifdefined\hypertarget\hypertarget{sourcepointer-4}{}\fi
\noindent\textit{Source reference.} Original reference 0040: sources/stacks-chow/tags.txt, source lines 156--156. Cross-chapter standard prerequisite cited by the author inside the counted span; it is used as a pointer and its own proof is not part of this item.\par\medskip
\ifdefined\hypertarget\hypertarget{sourcepointer-5}{}\fi
\noindent\textit{Source reference.} Original lemma 044V: sources/stacks-chow/tags.txt, source lines 5414--5414. Cross-chapter standard prerequisite cited by the author inside the counted span; it is used as a pointer and its own proof is not part of this item.\par\medskip
\ifdefined\hypertarget\hypertarget{sourcepointer-6}{}\fi
\noindent\textit{Source reference.} Original lemma 04Y4: sources/stacks-chow/tags.txt, source lines 6417--6417. Cross-chapter standard prerequisite cited by the author inside the counted span; it is used as a pointer and its own proof is not part of this item.\par\medskip
\begin{lemma}
\label{lemma-quotient-stack-2-cartesian}
Assume $B \to S$, $(U, R, s, t, c)$, and $\pi : \mathcal{S}_U \to [U/R]$
are as in
Lemma \ref{lemma-quotient-stack-arrows}.
The $2$-commutative square
$$
\xymatrix{
\mathcal{S}_R \ar[r]_s \ar[d]_t & \mathcal{S}_U \ar[d]^\pi \\
\mathcal{S}_U \ar[r]^-\pi & [U/R]
}
$$
of
Lemma \ref{lemma-quotient-stack-2-arrow}
is a $2$-fibre product of stacks in groupoids of $(\Sch/S)_{fppf}$.
\end{lemma}

\begin{proof}
According to
Stacks, Lemma \ref{stacks-lemma-2-product-stacks-in-groupoids}
the lemma makes sense. It also tells us that we have to show that
the functor
$$
\mathcal{S}_R \longrightarrow \mathcal{S}_U \times_{[U/R]} \mathcal{S}_U
$$
which maps $r : T \to R$ to $(T, t(r), s(r), \alpha(r))$ is an equivalence,
where the right hand side is the $2$-fibre product as described in
Categories, Lemma \ref{categories-lemma-2-product-categories-over-C}.
This is, after spelling out the definitions, exactly the content of
Lemma \ref{lemma-quotient-stack-morphisms}. (Alternative proof: Work out
the meaning of
Lemma \ref{lemma-cartesian-square-of-morphism}
in this situation will give you the result also.)
\end{proof}

\end{document}
